%% Derived from the official ACM sigconf sample kept in template/sigconf.tex.
%% Chinese internal draft: translate the body to English before submission.
\documentclass[sigconf,nonacm]{acmart}

\usepackage{amsmath,amssymb,booktabs,array,multirow}
\usepackage{xeCJK}
\setCJKmainfont{FandolSong-Regular}
\setCJKsansfont{FandolHei-Regular}
\setCJKmonofont{FandolFang-Regular}

\settopmatter{printacmref=false,printccs=false,printfolios=true}
\renewcommand\footnotetextcopyrightpermission[1]{}
\pagestyle{plain}

\newcommand{\method}{\textsc{SNSD-v5.1}}
\newcommand{\phased}{\textsc{Phased}}
\newcommand{\fullmethod}{\textsc{Full}}
\newcommand{\oneshot}{\textsc{OneShot}}
\newcommand{\selectiononly}{\textsc{SelectionOnly}}

\begin{document}

\title{SNSD-v5.1：面向星地链路调度的语义保护三图联合协商}

\author{Anonymous Authors}
\affiliation{%
  \institution{Anonymous Institution}
  \city{Anonymous}
  \country{Anonymous}}

\begin{abstract}
分组式多智能体调度会切断跨组冲突和收益替换关系，标量消息或搜索锚点难以恢复丢失的
组合结构。本文提出\method，耦合私有子图、动态边界图和公共全局图：局部图保留真实完整
候选，边界图组织跨智能体 delete--replace--refill 动作，全局图执行权威验证。方法从同一
不可变状态配对生成零语义和处理分支，仅允许严格改善且被接纳的语义响应影响未来候选域。
在完整 C3、6 智能体和每次 288 次相同 worker 调用下，结构化联盟动作在有/无 learned
semantics 时分别取得 $20{,}214.1$ 和 $16{,}424.2$ 的平均配对增益，均为 16/16 seeds
正向；持续分阶段通信相对零语义和一次性通信分别提高 $5{,}033.6$ 和 $4{,}810.7$。
结论限定于 C3 机制验证，不声称全局最优或跨实例泛化。
\end{abstract}

\keywords{multi-agent systems, satellite scheduling, graph coordination}

\maketitle

\section{方法概览}

本文关注的瓶颈并非单个智能体能否继续增强局部求解，而是强初始解附近如何传播和共同
可行化多个智能体的私有改进。分组已经第一次压缩了全局组合空间；若通信再把局部状态压成
一个标量或池化嵌入，就难以重构跨边界的一对多替换和回填。因此，\method 将信息流固定为
三个不对称阶段：
\begin{equation}
\text{局部子图生成响应}
\rightarrow
\text{边界图协商交换}
\rightarrow
\text{全局图验证并提交}.
\end{equation}

方法包含两个核心贡献。第一，结构保持的三图原始--对偶执行器把通信直接转换为完整的
跨智能体 delete--replace--refill 响应，而不只修改下轮搜索中心。第二，同状态
$C_0/C_A$ 配对保护阻止未接纳的语义响应污染未来搜索域，并保证已提交目标值单调不下降。
论文主配置 \phased（代码名 \texttt{AlternatingConflictBenefit}）进一步把冲突处理和收益
恢复按步骤交替消费。

\section{问题定义}

在固定规划窗口中，每个可执行通信机会对应顶点 $v\in V$，非负权重 $w_v$ 为链路持续
时间。当两个机会因为卫星或地面站资源冲突而不能同时调度时，在二者之间建立无向边。
问题写为最大权独立集：
\begin{equation}
\begin{aligned}
\max_{x\in\{0,1\}^{|V|}}\quad &F(x)=\sum_{v\in V}w_vx_v,\\
\text{s.t.}\quad &x_u+x_v\leq 1,\quad \forall (u,v)\in E.
\end{aligned}
\label{eq:mwis}
\end{equation}

按链路开始时间把顶点划分为 $K$ 个互不相交的分组，
$V=\dot\bigcup_{i=1}^{K}V_i$。智能体 $i$ 持有局部诱导图
$G_i=(V_i,E_i)$，其中
\begin{equation}
E_i=\{(u,v)\in E:u,v\in V_i\},\qquad
E_{\mathrm{cut}}=E\setminus\bigcup_iE_i.
\end{equation}
完整 C3 图包含 69,923 个顶点、979,960 条冲突边和 8,980 条切边，正式实现使用 6 个
时间智能体。所有方法使用相同 local worker；worker 返回实际实现的局部独立集，而不是被
假设为局部最优。多智能体协议负责选择局部状态、共享边界结构、组合真实局部记录并修复
跨组不兼容。

\section{三个耦合图}

\subsection{私有局部子图}

智能体不只保留当前锚点，而是维护少量 realized elites。每条 elite 包含 worker 实际返回的
完整顶点集、精确局部收益和边界签名。完整局部集合保留了删除边界顶点之后构造替代响应所需
的高维结构；其他智能体无需得到整个私有搜索空间，只接收真实记录及确定性的同组邻域扩展。

\subsection{动态边界交互图}

边界图由真实切边、当前 elites 暴露的边界顶点和响应候选动态构成。它传递明确的顶点身份、
冲突端点、候选角色、删除成本和潜在回填收益。仍存在切边冲突的智能体形成活跃交互图，其
大小为 2 或 3 的连通分量组成临时联盟。联盟共同决定删除哪些边界端点，受影响智能体再从
各自私有子图中执行替换和回填。

\subsection{公共全局图}

完整冲突图提供局部图和边界摘要无法安全近似的长距离冲突、真实目标和因子残差。它负责
生成压力/冲突/机会公共信号，并对每个完整响应执行权威独立性检查。全局图是公共事实层和
最终安全屏障，不替代私有搜索，也不读取未查询的未来 worker 结果。

\section{结构保持的原始--对偶联合动作}

\subsection{带身份的边界包}

后期协商时，智能体 $i$ 保留最多 $T$ 个价值最高且互异的真实记录，构造
\begin{equation}
P_i=(h_i,\mathcal E_i,B_i,C_i,\Phi_i,E_i^{\mathrm{cut}},R_i),
\end{equation}
其中 $h_i$ 是锚点哈希，$\mathcal E_i$ 是 elite ID，$B_i$ 是已选边界顶点，$C_i$ 是
有限真实候选，$\Phi_i$ 是候选特征，$E_i^{\mathrm{cut}}$ 是相关切边，$R_i$ 是类型化关系
计数。每个候选 $v\in C_i$ 的特征为
\begin{equation}
\phi_i(v)=(f_v,w_v,\widehat\Delta_v,d_v,\widehat r_v,u_v,\rho_v),
\end{equation}
分别表示 elite 选择频率、原始收益、经验局部边际、删除成本、最佳回填收益、不确定性和
\texttt{keep}/\texttt{add-or-refill} 操作角色。所有 ID 必须来自已实现记录或确定性的同组
一跳/两跳扩展，不能由未来宏轮结果制造。

\subsection{三种语义通道}

图消息管线为 $v\in V_i$ 输出压力 $p_i(v)$、冲突 $c_i(v)$ 和机会 $o_i(v)$，并在每个
智能体内按最大绝对值归一化。完整语义信用为
\begin{equation}
q_i^{\mathrm{full}}(v)=
\operatorname{norm}_i\left[p_i(v)+c_i(v)-\tfrac12|o_i(v)|\right].
\end{equation}
零语义控制令 $q_i(v)=0$。分阶段配置分别使用
\begin{equation}
\begin{split}
q_i^{\mathrm{conf}}(v)&=\operatorname{norm}_i[p_i(v)+c_i(v)],\\
q_i^{\mathrm{benefit}}(v)&=\operatorname{norm}_i[p_i(v)-o_i(v)].
\end{split}
\end{equation}
Shuffled 控制保持消息数值集合、通信量和计算预算不变，但将通道值重新分配给错误的
智能体/顶点 ID，以破坏语义对应关系。

\subsection{切边价格和联合删除}

令 $S^{\mathrm{priv}}$ 为当前私有 elites 并集。在微循环 $m$ 中，每条活跃切边
$e=(u,v)$ 的非负价格按下式更新：
\begin{equation}
\lambda_e^{m+1}=\max\left\{0,\lambda_e^m+
\frac{0.20}{\sqrt{m+1}}\max\left(1,\frac{w_u+w_v}{2}\right)\right\}.
\end{equation}
顶点价格为 $\Lambda_v=\sum_{e\ni v}\lambda_e$。两个冲突端点的保留分数为
\begin{equation}
\begin{aligned}
k_i(v)&=w_v-\min\{w_v,\widehat r_v\}
+0.35s_{uv}q_i(v)-\Lambda_v,\\
s_{uv}&=\max\left(1,\frac{w_u+w_v}{2}\right).
\end{aligned}
\label{eq:keep}
\end{equation}
执行器暂时删除分数较低的端点，并以固定收益/ID 规则处理平局。删除在联盟连通分量中分批
执行，直到所有活跃切边冲突消失。式~\eqref{eq:keep} 是离散非凸协调启发式，本文不作
ADMM 或渐近收敛声明。

\subsection{Delete--replace--refill 闭包}

对每个受删除影响的智能体，从被删顶点的同组一跳/两跳邻域以及保留 elites 中构造回填池。
候选排序为
\begin{equation}
r_i(v)=\frac{w_v}{1+\log(1+\deg(v))}+0.20w_vq_i(v).
\end{equation}
执行器贪心加入与当前全局选择兼容的候选，再尝试确定性 one-for-many exchange：只有新顶点
的冲突全部位于同一本地分组，且其收益严格大于被替换顶点总收益时才执行交换。输出是所有
智能体的完整响应，而不是新锚点；完整图随后验证可行性并计算真实目标。

\section{同状态语义保护}

\subsection{配对 \texorpdfstring{$C_0/C_A$}{C0/CA} 分支}

一个当前未被选择的语义响应如果仍进入后续 elites，就会改变未来候选域，造成语义域污染。
为此，每个后期决策先冻结同一不可变状态，包括私有 elites、当前记录、联合记录、完整图和
消息通道；随后配对生成零语义响应 $C_0$ 与注册处理响应 $C_A$。当处理本身为零时，
$C_A\equiv C_0$，第二次执行被去重。否则两条分支只在语义信用上不同。

\subsection{反事实微响应}

第 0 个微循环复现原始结构化响应，后续微循环使用上一响应的删除和回填计数构造不同方案：
\begin{equation}
\widetilde q_i^m(v)=q_i(v)+5n_v^{\mathrm{del}}-3n_v^{\mathrm{refill}}.
\end{equation}
曾被删除的端点获得保留奖励，曾被回填的顶点受到再次选择惩罚。所有结构互异的响应都经
全图验证，分支返回真实目标最高者。因此多微循环形成反事实删除/回填组合，而不是重复定价
同一个动作。

\subsection{严格语义门和单调提交}

后期决策 $t$ 的候选包括当前调度 $Y_t$、因子/链组合 $J_t$、无消息修复 $R_t^0$、
消息条件修复 $R_t^m$、零语义联盟响应 $C_t^0$ 和处理响应 $C_t^A$。处理分支只有满足
\begin{equation}
F(C_t^A)>\max\{F(J_t),F(R_t^0),F(R_t^m),F(C_t^0)\}
\label{eq:gate}
\end{equation}
才具有资格。当联合动作因子 $B=1$ 时，有限候选并集为
\begin{equation}
\mathcal U_t=\{Y_t,J_t,R_t^0,R_t^m,C_t^0\}
\cup\{C_t^A:\text{式~\eqref{eq:gate}成立}\}.
\end{equation}
当 $B=0$ 时，$C_t^0$ 和 $C_t^A$ 仍可用于配对诊断，但不进入提交并集。最终选择为
\begin{equation}
Y_{t+1}=\arg\max_{Y\in\mathcal U_t,\,Y\text{ 全图可行}}F(Y).
\end{equation}
因为 $Y_t\in\mathcal U_t$，所以 $F(Y_{t+1})\ge F(Y_t)$。该性质只保证提交序列单调，
不保证式~\eqref{eq:mwis} 的全局最优。

未来候选域采用独立安全规则：$C_t^0$ 在 $B=1$ 时作为受保护联盟基线；$C_t^A$ 只有通过
严格门且最终被选为当前锚点时才能进入下一域。未接纳 $C_t^A$ 的任何组件均不注入。

\section{持续冲突--收益分阶段协议}

\phased 在后期 $C_A$ 分支中按源步骤交替使用 $q^{\mathrm{conf}}$ 和
$q^{\mathrm{benefit}}$。零基实现的首个后期源步骤是 5，因此序列以冲突阶段开始，再进入
收益阶段。冲突阶段回答“为了共同可行，谁应让步”，收益阶段回答“让步后如何恢复价值”。
\fullmethod 在每个后期状态同时使用所有通道；\oneshot 只在第一个后期状态使用真实语义，
此后使用 $C_0$；Shuffled 保持载荷和频次但破坏语义归属。

需要区分同状态效应 $F(C_A)-F(C_0)$ 与端到端 A 因子。前者只衡量当前联合执行器的即时
语义干预，后者还包含早期消息对 worker 提案和 realized candidate trajectory 的影响。

\section{完整执行流程}

每个宏轮首先由当前全局调度初始化各智能体的私有 elite 和记录域。在每个步骤中，6 个
智能体并行调用 local worker 各一次，更新真实局部记录，并由因子/链协调器构造 $J_t$ 和
边界残差。前 6 步只开发私有搜索；后 6 步依次构造 $R_t^0$、可用的 $R_t^m$，冻结状态，
配对执行 $C_t^0/C_t^A$，应用式~\eqref{eq:gate}，最后在完整图上选择目标最高者。

正式协议固定 4 个宏轮、每轮 12 步和 6 个智能体，因此每次运行使用
\begin{equation}
4\times 12\times 6=288
\end{equation}
次付费 worker 调用。结构化分支、微循环和候选验证增加 0 次 worker 调用，且不因连续无提升
而提前终止。每个宏轮有 6 个后期决策，整次运行共 24 个后期联合协商状态。

\section{通信和计算开销}

令 $T_i\le T$ 为智能体 $i$ 的 elite 数，$c_i\le4T$ 为 packet 候选数，$b_i$ 为已选
边界 ID 数，$r_i$ 为类型化关系数。实现按
\begin{equation}
L_i=8c_i+3r_i+T_i+b_i+6
\end{equation}
个数值槽位核算一个 packet，ID 和哈希保守按 8 字节计。含 $M$ 个微循环的一条分支为
\begin{equation}
L_{\mathrm{branch}}=M\sum_iL_i+2\sum_{m=0}^{M-2}
\left(|D_m\cap B|+|H_m|\right),
\end{equation}
报告字节数为 $8L_{\mathrm{branch}}$。这是一致的载荷核算约定，不代表某个网络序列化器的
实测开销。

packet 提取复杂度为
$O(\sum_iT_i|B_i|+\sum_ic_i^2+H_{\mathrm{packet}})$；单个微循环为
$O(|E_{\mathrm{cut}}(S^{\mathrm{priv}})|+H_{\mathrm{refill}}+|V|+|E|)$，其中最后一项
是权威全图验证。实现把回填池限制为 192 个顶点，把交换短名单限制为 64 个顶点，最终最多
验证 6 个命名候选。冻结实现使用 $T\in\{4,6\}$、$M\in\{3,4\}$。

\section{Learned 和解析模式}

v5.1 执行器本身没有端到端重新训练。Learned 运行严格加载验证选择的 TriGraph checkpoint
并冻结参数；GNN/消息模型负责继承的提案和边界通道，结构化执行器将其转为联合响应。该模型
含 456,075 个参数，但没有使用 v5.1 特有的 $C_A-C_0$ 或 coalition-delta 监督。

解析 Oracle 移除 checkpoint，通过同一接口使用解析压力/排序脚手架，只诊断机制能否利用
信息充分的冲突和机会信号。它不知道未查询的未来候选，不是全局 MWIS 上界，也不是 Learned
运行时教师；两类结果必须分开报告。

\section{正式实验协议}

所有实验均使用完整 C3 而非 smoke。Original Learned 使用 seeds
$\{11,22,33,44,55,66,77,88\}$；Fresh Learned 使用在矩阵执行前固定的 seeds
$\{101,202,303,404,505,606,707,808\}$；Oracle 机制诊断使用 Original 8 seeds。
每套矩阵包含表~\ref{tab:methods} 的 8 个配置，合计 64 runs。

\begin{table}[t]
\caption{冻结消融矩阵。A 为真实语义处理，B 为结构化联盟响应进入选择和受保护域。}
\label{tab:methods}
\small
\begin{tabular}{lccc}
\toprule
配置 & A & B & 作用 \\
\midrule
Full & 1 & 1 & 同时语义 + 联盟 \\
FrozenAnchorComm & 1 & 0 & 语义、无联盟注入 \\
SelectionOnly & 0 & 1 & 零语义 + 联盟 \\
FrozenAnchorSelectionOnly & 0 & 0 & 因子基线 \\
AlternatingConflictBenefit & 1 & 1 & 分阶段持续语义 \\
OneShotComm & 部分 & 1 & 仅一次真实语义 \\
ShuffledComm & 置乱 & 1 & 错误语义归属 \\
NoPeerFinalShield & 0 & 0 & 无 peer、仅安全屏障 \\
\bottomrule
\end{tabular}
\end{table}

三套矩阵共 192 runs，每次 288 次 worker 调用，总计 55,296 次。两套 Learned 矩阵均严格
加载同一个 456,075 参数 checkpoint；其文件 SHA-256 的完整值随复现包清单发布。
128/128 Learned runs 均满足严格键匹配且 \texttt{oracle\_mode=false}。三套矩阵的缺失/意外/重复
cell、预算违规、额外 worker 调用、语义日程违规和未接纳 $C_A$ 污染均为 0。

\section{Learned 结果}

增益定义为 $F(Y_{\mathrm{final}})-F(Y_{\mathrm{initial}})$，表中报告 seed 间均值和样本
标准差。Pooled 16 仅用于描述配对效应，不能解释为 16 个独立问题实例。

\begin{table}[t]
\caption{完整 C3 上的 Learned 最终增益（mean $\pm$ sample SD）。}
\label{tab:learned-main}
\centering
\scriptsize
\setlength{\tabcolsep}{2.2pt}
\begin{tabular}{lrrr}
\toprule
方法 & Original 8 & Fresh 8 & Pooled 16 \\
\midrule
\textbf{Phased} & \textbf{38,805.5$\pm$5,752.8} &
\textbf{32,963.8$\pm$5,022.8} & \textbf{35,884.6$\pm$6,026.4} \\
Full & 36,630.4$\pm$7,956.1 & 31,149.0$\pm$5,115.7 & 33,889.7$\pm$7,054.4 \\
Shuffled & 35,894.3$\pm$8,015.7 & 32,725.3$\pm$6,467.9 & 34,309.8$\pm$7,223.9 \\
OneShot & 34,777.9$\pm$7,994.1 & 27,370.0$\pm$5,702.9 & 31,073.9$\pm$7,722.3 \\
SelectionOnly & 34,133.1$\pm$5,945.2 & 27,569.0$\pm$7,530.3 & 30,851.1$\pm$7,378.8 \\
FrozenComm & 15,106.4$\pm$6,545.1 & 12,244.9$\pm$4,849.4 & 13,675.6$\pm$5,757.6 \\
A0B0 & 16,416.1$\pm$4,987.8 & 12,437.6$\pm$5,277.3 & 14,426.9$\pm$5,369.1 \\
NoPeer & 16,423.0$\pm$4,983.1 & 12,437.6$\pm$5,277.3 & 14,430.3$\pm$5,368.4 \\
\bottomrule
\end{tabular}
\end{table}

\phased 在 Original、Fresh 和 pooled 三个视图中均取得最高平均增益，其平均最终目标分别为
6,121,118.6、6,116,224.0 和 6,118,671.3。不过 Fresh 上 \phased 只比 Shuffled 高
238.5，因而只能称为平均表现最好的冻结配置，不能声称准确 learned 语义稳定击败置乱语义。

\subsection{结构化联盟动作是最稳定的主效应}

\begin{table}[t]
\caption{Pooled 16-seed 因子 B 的配对效应。}
\label{tab:factor-b}
\centering
\small
\begin{tabular}{lrrr}
\toprule
效应 & Mean $\pm$ SD & 正向 & Sign $p$ \\
\midrule
$B\mid A$ & \textbf{20,214.1 $\pm$ 5,685.9} & \textbf{16/16} & $3.05\!\times\!10^{-5}$ \\
$B\mid\neg A$ & \textbf{16,424.2 $\pm$ 5,095.9} & \textbf{16/16} & $3.05\!\times\!10^{-5}$ \\
\bottomrule
\end{tabular}
\end{table}

无论有无 learned semantics，结构化联盟响应进入候选选择和后续受保护域后都在 16/16 seeds
产生大幅正向提升。这是当前最强证据：主要收益来自 delete--replace--refill 联合动作实质
扩展了全局可行候选域，而非某个消息排序器的偶然结果。

\subsection{持续分阶段通信优于零语义和一次通信}

\begin{table}[t]
\caption{Pooled 16-seed 通信协议配对差。}
\label{tab:protocol-effects}
\centering
\small
\begin{tabular}{lrrr}
\toprule
对比 & Mean $\pm$ SD & 正向 & Sign $p$ \\
\midrule
Phased $-$ SelectionOnly & \textbf{5,033.6 $\pm$ 5,072.7} & \textbf{14/16} & .0042 \\
Phased $-$ OneShot & \textbf{4,810.7 $\pm$ 4,742.8} & \textbf{14/16} & .0042 \\
Phased $-$ Full & 1,994.9 $\pm$ 4,759.3 & 10/16 & .4545 \\
Phased $-$ Shuffled & 1,574.9 $\pm$ 5,799.1 & 12/16 & .0768 \\
\bottomrule
\end{tabular}
\end{table}

前两项支持持续冲突--收益通信的附加价值；对应的双侧 sign-flip mean $p$ 分别为 .0017 和
.0019。后两项只显示平均方向为正，不能支持普遍或统计显著支配。

语义因子在联盟存在时为 $A\mid B=3,038.6\pm5,773.6$，12/16 正向，精确 sign
$p=.0768$，属于边缘补充证据；脱离联盟时为
$A\mid\neg B=-751.3\pm5,489.2$。因此更稳妥的解释是：语义信息需要与可执行的联合动作
耦合，而不是独立地产生稳定主效应。

\section{Oracle 机制诊断与局限}

Original 8 seeds 中，Oracle Full 的增益为 $35,589.1\pm4,433.3$，OneShot 为
$36,860.0\pm3,965.6$；$B\mid A=17,428.5\pm5,358.2$ 和
$B\mid\neg A=17,994.8\pm4,135.2$ 均为 8/8 正向。该矩阵再次验证联合动作机制能利用信息
充分的冲突/机会信号，但 Oracle 不是全局上界，也不能和 Learned 配置排名混为一谈。

局限为：单一 C3；未与 KP-NLNS/CPLEX 作同协议比较；Phased 对 Full/Shuffled 的优势
不均匀；checkpoint 未以 v5.1 联盟差值重训。因此本文只主张联合机制及配对效应，不主张
跨实例泛化、绝对 SOTA 或全局最优。


\end{document}
